% The appendix deliberately uses only the packages already loaded by
% compose_arxiv.sty.  Its ruled pseudocode follows the compact visual rhythm of
% the reference manuscript without introducing an algorithm package.
\newcounter{composeappendixalgorithm}
\newenvironment{composeappendixalgorithm}[1]{%
  \refstepcounter{composeappendixalgorithm}%
  \par\medskip\noindent
  \begin{minipage}{\linewidth}
  \hrule\vspace{0.28em}
  \textbf{Algorithm \thecomposeappendixalgorithm}\quad #1\par
  \vspace{0.22em}\hrule\vspace{0.38em}
}{%
  \vspace{0.18em}\hrule
  \end{minipage}
  \par\medskip
}

\newtheorem{composeappendixproposition}{Proposition}[section]
\newtheorem{composeappendixlemma}[composeappendixproposition]{Lemma}

\section{Algorithms}
\label{sec:app-compose-algorithms}

This appendix states the finite-support procedures corresponding to the
definitions in Sections~\ref{sec:rgm}--\ref{sec:finite-control}.  The
procedures separate three objects that must not be conflated: the learned law
over executable marks, its canonical molecular pushforward, and a controller
over productive molecular successors.  They also separate fixed-budget
editing from timed de novo simulation.  Only the former is used below.

\subsection{Training losses on marks and molecular successors}
\label{sec:app-compose-training}

Algorithm~\ref{alg:app-compose-training} constructs one training-row loss.
The objective mode is part of the checkpoint contract; selecting a different
mode changes the trained stochastic object.  A productive teacher row carries
an executable mark $a^\star$, its stored canonical successor $y^\star$, a
teacher rate $r>0$, and an importance weight $w$.  A terminal row has $r=0$.
The productive-identity mode excludes terminal rows from its normalized batch
average, whereas the two rate objectives retain their appropriate positive
hazard term.

\begin{composeappendixalgorithm}{Mark- or successor-level training-row loss}
\label{alg:app-compose-training}
\textbf{Require:} state $x$, progress $t$, exact legal fiber $\A(x)$,
executor $\T$, canonical identity $c$, prediction
$(\Lambda_\theta,p_\theta)$, teacher row
$(a^\star,y^\star,r,w)$, and a declared objective mode.
\begin{enumerate}[
  leftmargin=2.45em,
  label=\arabic*:,
  topsep=0.25em,
  itemsep=0.12em,
  parsep=0pt
]
  \item Enumerate the complete declared fiber $\A(x)$ and verify
  $\T(x,a)\in\X$ for every $a\in\A(x)$.
  \item Normalize the factorized scores so that
  $\sum_{a\in\A(x)}p_\theta(a\mid x,t)=1$, and compute
  $\Lambda_\theta(x,t)\geq0$.
  \item Execute and canonicalize every mark, then group its probability by
  $y=c(\T(x,a))$ to obtain $\bar p_\theta$, $v_\theta$, and $z_\theta$
  using Algorithm~\ref{alg:app-compose-kernel}.
  \item If $r>0$, require $a^\star\in\A(x)$,
  $c(\T(x,a^\star))=y^\star\neq c(x)$, and positive model mass on the
  teacher coordinate required by the selected objective.
  \item For \emph{selected-mark GM}, return
  \[
    w\!\left[
      \Lambda_\theta
      -r\{\log\Lambda_\theta+\log p_\theta(a^\star\mid x,t)\}
    \right],
  \]
  with the $r=0$ row equal to $w\Lambda_\theta$.
  \item For \emph{productive identity}, if $r>0$ return
  $-w\log\{\bar p_\theta(y^\star\mid x,t)/z_\theta(x,t)\}$; otherwise
  mark the row as excluded from this loss.
  \item For \emph{successor GM}, return
  \[
    w\!\left[
      \Lambda_\theta z_\theta
      -r\{\log\Lambda_\theta+\log\bar p_\theta(y^\star\mid x,t)\}
    \right],
  \]
  with the $r=0$ row equal to
  $w\Lambda_\theta z_\theta$.
\end{enumerate}
\end{composeappendixalgorithm}

The selected-mark expression is
Equation~\eqref{eq:mark-gm-loss}; productive identity is the rowwise term of
Equation~\eqref{eq:productive-identity-loss}; and successor GM is
Equation~\eqref{eq:successor-gm-loss}.  No arbitrary subset of candidates can
replace $\A(x)$ without redefining the relevant normalizer and therefore the
model law.

\subsection{Canonical grouping and the productive editing kernel}
\label{sec:app-compose-kernel}

\begin{composeappendixalgorithm}{Canonical grouping and productive-kernel construction}
\label{alg:app-compose-kernel}
\textbf{Require:} executor state $x$, time $t$, finite legal fiber $\A(x)$,
executor $\T$, canonical identity $c$, total marked hazard
$\Lambda_\theta$, and normalized mark probabilities $p_\theta$.
\begin{enumerate}[
  leftmargin=2.45em,
  label=\arabic*:,
  topsep=0.25em,
  itemsep=0.12em,
  parsep=0pt
]
  \item Initialize a zero-valued map $M$ from canonical identities to
  nonnegative mass.
  \item For every $a\in\A(x)$, execute $x_a\leftarrow\T(x,a)$, set
  $y_a\leftarrow c(x_a)$, and accumulate
  $M[y_a]\leftarrow M[y_a]+p_\theta(a\mid x,t)$.
  \item Set $\bar p_\theta(y\mid x,t)\leftarrow M[y]$,
  $v_\theta(x,t)\leftarrow M[c(x)]$, and
  $z_\theta(x,t)\leftarrow1-v_\theta(x,t)$.
  \item For every $y\neq c(x)$, set the off-diagonal molecular rate
  $Q_{\theta,t}(x,y)\leftarrow
  \Lambda_\theta(x,t)\bar p_\theta(y\mid x,t)$.
  \item Set
  $\Lambda_\theta^{\mathrm{prod}}(x,t)\leftarrow
  \Lambda_\theta(x,t)z_\theta(x,t)$.
  \item If $z_\theta(x,t)>0$, return the productive editing row
  $P_{\theta,t}^{\mathrm{edit}}(y\mid x)
  \leftarrow\bar p_\theta(y\mid x,t)/z_\theta(x,t)$ for
  $y\neq c(x)$.
  \item If $z_\theta(x,t)=0$, return \emph{terminal under the productive
  editing convention}; do not fabricate a self-transition.
\end{enumerate}
\end{composeappendixalgorithm}

Algorithm~\ref{alg:app-compose-kernel} is the scientific dictionary
definition.  A segmented reduction may compute the same sums more efficiently,
but it is not a different kernel.  The canonical map includes the distinguished
null identity; the null state is still not a molecule.

\subsection{Fixed-budget editing}
\label{sec:app-compose-fixed-budget}

Let $t_b$ be the registered model progress used when $b$ productive edits
remain.  If model progress changes with the budget, then
$P_b=P_{\theta,t_b}^{\mathrm{edit}}$ is budget indexed.  Sampling a canonical
successor does not by itself provide an executable action, so the algorithm
also samples a conditional lift inside the chosen successor fiber.

\begin{composeappendixalgorithm}{Fixed-budget productive editing}
\label{alg:app-compose-fixed-budget}
\textbf{Require:} exact source state $x_0$, maximum productive budget $K$,
registered progress values $\{t_b\}_{b=1}^K$, and optionally a
successor-level controlled row $R_b(\cdot\mid x)$; by default $R_b=P_b$.
\begin{enumerate}[
  leftmargin=2.45em,
  label=\arabic*:,
  topsep=0.25em,
  itemsep=0.12em,
  parsep=0pt
]
  \item Set $x\leftarrow x_0$.
  \item For $b=K,K-1,\ldots,1$, construct
  $\bar p_{\theta,t_b}$ and $P_b$ by
  Algorithm~\ref{alg:app-compose-kernel}.
  \item If $x$ is terminal under the productive convention, return the
  current trajectory and record that fewer than $K$ jumps were available.
  \item Sample a canonical successor $y\sim R_b(\cdot\mid x)$.
  Require $\operatorname{supp}R_b(x,\cdot)
  \subseteq\operatorname{supp}P_b(x,\cdot)$.
  \item Sample a mark inside its fiber,
  \[
    a\sim
    \frac{p_\theta(a\mid x,t_b)}
         {\bar p_\theta(y\mid x,t_b)}
    \quad\text{over}\quad
    \{a\in\A(x):c(\T(x,a))=y\}.
  \]
  \item Form $x^+\leftarrow\T(x,a)$ only after rechecking legality and
  $c(x^+)=y$.  If the process uses a fixed canonical section, set
  $x\leftarrow s(y)$; otherwise retain $x\leftarrow x^+$ only under the
  representative-consistency condition in Assumption~A3.
  \item Return the exact committed-state trajectory and its canonical
  identities.
\end{enumerate}
\end{composeappendixalgorithm}

The hazard magnitude $\Lambda_\theta$ is absent from this procedure.  Timed
Gillespie simulation is a separate inference regime.

\subsection{Exact backward values and dynamic continuation}
\label{sec:app-compose-control-algorithms}

\begin{composeappendixalgorithm}{Exact finite-horizon backward values and Doob rows}
\label{alg:app-compose-doob}
\textbf{Require:} an exactly enumerable state set, row-stochastic
budget-indexed kernels $\{P_b\}_{b=1}^K$, and a finite nonnegative terminal
desirability $g$.
\begin{enumerate}[
  leftmargin=2.45em,
  label=\arabic*:,
  topsep=0.25em,
  itemsep=0.12em,
  parsep=0pt
]
  \item Set $h_0(x)\leftarrow g(x)$ for every state $x$.
  \item For $b=1,\ldots,K$, compute
  $h_b(x)\leftarrow\sum_yP_b(y\mid x)h_{b-1}(y)$ for every $x$.
  \item If $h_b(x)=0$, mark $(b,x)$ \emph{unreachable for this target and
  budget}; leave its controlled row undefined.
  \item If $h_b(x)>0$, set
  \[
    P_b^g(y\mid x)
    \leftarrow
    P_b(y\mid x)\frac{h_{b-1}(y)}{h_b(x)}
  \]
  for every $y$ in the base row.
  \item Verify row normalization and
  $\operatorname{supp}P_b^g(x,\cdot)
  \subseteq\operatorname{supp}P_b(x,\cdot)$ on every defined row.
  \item Return $\{h_b\}_{b=0}^K$, all defined controlled rows, and the
  unreachable set at every remaining budget.
\end{enumerate}
\end{composeappendixalgorithm}

\begin{composeappendixalgorithm}{Exact continuation after an objective change}
\label{alg:app-compose-dynamic}
\textbf{Require:} current state $x$, remaining budget $b$, future base kernels
$P_1,\ldots,P_b$, and a newly declared terminal desirability $g'$.
\begin{enumerate}[
  leftmargin=2.45em,
  label=\arabic*:,
  topsep=0.25em,
  itemsep=0.12em,
  parsep=0pt
]
  \item Discard the old objective for future control, but retain the current
  state and all context required to make the future base law Markov.
  \item Run Algorithm~\ref{alg:app-compose-doob} with horizon $b$ and
  desirability $g'$.
  \item If $h_b^{g'}(x)=0$, report the new target unreachable from the current
  state under the remaining budget.
  \item Otherwise sample the next successor from
  $(P_b)^{g'}(\cdot\mid x)$ and continue with remaining-budget rows
  $(P_{b-1})^{g'},\ldots,(P_1)^{g'}$.
  \item Report this law as the exact $g'$-tilted \emph{continuation from
  $(b,x)$}, not as the counterfactual law obtained by using $g'$ from the
  original source.
\end{enumerate}
\end{composeappendixalgorithm}

\section{Proofs and formal scope}
\label{sec:app-compose-proofs}

\subsection{Assumptions and the quotient-state boundary}
\label{sec:app-compose-assumptions}

The one-step pushforward requires only a finite legal fiber, deterministic
execution, and a normalized marked law.  A multi-step Markov process on
canonical identities requires one additional interface condition.  We state
it explicitly.

\begin{enumerate}[
  leftmargin=1.65em,
  label=\textbf{A\arabic*.},
  topsep=0.25em,
  itemsep=0.22em
]
  \item \textbf{Finite executable fibers.}  For each executor state $x$,
  $\A(x)$ is finite; every enumerated mark executes deterministically to a
  state in $\X$.
  \item \textbf{Finite marked rate.}  The mark law is normalized and
  $\Lambda_\theta(x,t)<\infty$.
  \item \textbf{Quotient consistency.}  Either every committed canonical
  identity $y$ is converted to a fixed exact representative $s(y)$ before
  the next row is evaluated, or for any exact representatives $x,x'$ with
  $c(x)=c(x')$ and every $y\neq c(x)$,
  \begin{equation}
    \sum_{a:\,c(\T(x,a))=y}\lambda_\theta(a\mid x,t)
    =
    \sum_{a':\,c(\T(x',a'))=y}
    \lambda_\theta(a'\mid x',t).
    \label{eq:app-compose-lumpability}
  \end{equation}
  The analogous equality is required for any budget-indexed embedded rows.
  \item \textbf{Finite-horizon control.}  Each $P_b$ is row stochastic on a
  finite or countable state set; $g\geq0$ is finite; and all backward sums used
  below are finite.
  \item \textbf{Sufficient control state.}  Conditional on the current state,
  remaining budget, and declared control context, the future base law is
  independent of the earlier trajectory.  History-dependent constraints are
  included by augmenting the state or are outside the theorem.
\end{enumerate}

Assumption~A3 is not implied merely by observing that several marks share a
one-step canonical successor.  It is the precise extra antecedent needed to
call the multi-step canonical quotient Markov.  The one-step pushforward and
refinement identities below are pointwise without A3; every multi-step
molecular statement assumes a well-defined quotient kernel under A3.  This
distinction also prevents persistent-slot aliases with different future rows
from being silently merged.

\subsection{Pushforward generator and productive hazard}
\label{sec:app-compose-pushforward-proof}

\begin{composeappendixproposition}[Canonical pushforward and productive exit]
\label{prop:app-compose-pushforward}
Under A1--A2, for every bounded test function
$f$ on canonical identities,
\begin{equation}
  \bigl(\mathcal L_{\theta,t}(f\circ c)\bigr)(x)
  =
  \sum_{y\neq c(x)}
  \Lambda_\theta(x,t)\bar p_\theta(y\mid x,t)
  \{f(y)-f(c(x))\}.
  \label{eq:app-compose-pushed-generator}
\end{equation}
Hence the molecular off-diagonal rate is
$Q_{\theta,t}(x,y)=\Lambda_\theta\bar p_\theta(y\mid x,t)$,
its productive exit hazard is
$\Lambda_\theta^{\mathrm{prod}}=\Lambda_\theta(1-v_\theta)$,
and, when $z_\theta=1-v_\theta>0$, its embedded productive row is
$P_{\theta,t}^{\mathrm{edit}}(y\mid x)=
\bar p_\theta(y\mid x,t)/z_\theta$.  Under A3 these rows descend to a
well-defined Markov process on canonical identities.
\end{composeappendixproposition}

\paragraph{Proof.}
Apply the marked generator in Equation~\eqref{eq:marked-generator} to
$f\circ c$ and group marks by $y=c(\T(x,a))$:
\begin{align}
  \bigl(\mathcal L_{\theta,t}(f\circ c)\bigr)(x)
  &=
  \sum_{a\in\A(x)}
  \Lambda_\theta p_\theta(a\mid x,t)
  \{f(c(\T(x,a)))-f(c(x))\} \notag\\
  &=
  \sum_y
  \Lambda_\theta\bar p_\theta(y\mid x,t)
  \{f(y)-f(c(x))\}.
\end{align}
The $y=c(x)$ term is zero.  This proves
Equation~\eqref{eq:app-compose-pushed-generator} and identifies the
off-diagonal rates.  Summing them gives
\[
  \sum_{y\neq c(x)}Q_{\theta,t}(x,y)
  =
  \Lambda_\theta
  \sum_{y\neq c(x)}\bar p_\theta(y\mid x,t)
  =
  \Lambda_\theta(1-v_\theta)
  =
  \Lambda_\theta z_\theta.
\]
Conditioning the next molecular state change on a productive marked event
divides each off-diagonal rate by this exit hazard, yielding
$\bar p_\theta(y\mid x,t)/z_\theta$.  If $z_\theta=0$, the marked process may
still contain canonical self-events, but the canonical molecular state has no
exit and the productive embedded row is undefined.  Finally,
Equation~\eqref{eq:app-compose-lumpability}, or evaluation through a fixed
section $s$, makes each off-diagonal aggregate depend only on the canonical
source; this is exactly what is needed for the row to descend.
\hfill$\square$

\subsection{Within-fiber refinement invariance}
\label{sec:app-compose-refinement-proof}

\begin{composeappendixproposition}[Within-fiber invariance]
\label{prop:app-compose-refinement}
Consider two internal mark encodings with the same total hazard.  Suppose the
second refines the first only within canonical successor fibers: there is a
map $\rho:\A'(x)\to\A(x)$ such that
$c(\T'(x,a'))=c(\T(x,\rho(a')))$ and
\[
  \sum_{a':\,\rho(a')=a}p'_\theta(a'\mid x,t)
  =
  p_\theta(a\mid x,t)
  \quad\text{for every }a\in\A(x).
\]
Then $\bar p_\theta$, $v_\theta$, $z_\theta$, the molecular off-diagonal
rates, and the productive editing row are identical.  If A3 holds for both
encodings, every finite-horizon successor-level value and Doob row is also
identical.
\end{composeappendixproposition}

\paragraph{Proof.}
For any canonical successor $y$,
\begin{align}
  \bar p'_\theta(y\mid x,t)
  &=
  \sum_{a':\,c(\T'(x,a'))=y}p'_\theta(a'\mid x,t) \notag\\
  &=
  \sum_{a:\,c(\T(x,a))=y}
  \sum_{a':\,\rho(a')=a}p'_\theta(a'\mid x,t) \notag\\
  &=
  \sum_{a:\,c(\T(x,a))=y}p_\theta(a\mid x,t)
  =
  \bar p_\theta(y\mid x,t).
\end{align}
Taking $y=c(x)$ proves equality of $v_\theta$ and hence $z_\theta$.
Multiplication by the common total hazard proves equality of the
off-diagonal rates, and productive normalization proves equality of the
embedded rows whenever $z_\theta>0$.  For finite-horizon control, the base
rows are therefore equal.  The equality $h'_0=h_0=g$ and induction in
$h_b=P_bh_{b-1}$ give $h'_b=h_b$ for every $b$, after which the controlled-row
formula gives $(P'_b)^g=P_b^g$ on every reachable row.
\hfill$\square$

\begin{composeappendixproposition}[A mark-level power tilt is not refinement invariant]
\label{prop:app-compose-mark-negative}
For any positive power exponent other than one, a controller that powers
individual mark probabilities before canonical grouping is not invariant in
general to within-fiber refinement.
\end{composeappendixproposition}

\paragraph{Proof by counterexample.}
Let two canonical successors $y$ and $z$ each have base mass $1/2$.  In the
first encoding, one mark of mass $1/2$ realizes each successor.  In the second,
the $y$ mark is refined into four aliases of mass $1/8$, while the $z$ mark is
unchanged.  Both encodings therefore induce the same base successor law.
For an exponent $\beta>0$, the first encoding assigns equal powered weights
$2^{-\beta}$ to $y$ and $z$.  The refined encoding instead assigns powered
weights $4(1/8)^\beta$ and $(1/2)^\beta$.  Their ratio is
$2^{2-2\beta}$, which equals one if and only if $\beta=1$.  Thus every
nonidentity power can change the successor law under this refinement.  For
example, under a square tilt the first encoding assigns unnormalized
successor weights
\[
  \left(\frac12\right)^2=\frac14
  \quad\text{and}\quad
  \left(\frac12\right)^2=\frac14,
\]
so its normalized successor law remains $(1/2,1/2)$.  The refined encoding
assigns
\[
  4\left(\frac18\right)^2=\frac1{16}
  \quad\text{and}\quad
  \left(\frac12\right)^2=\frac14,
\]
so its normalized successor law is $(1/5,4/5)$.  The base molecular law was
unchanged, but the mark-level controller changed.  This proves the power-tilt
claim.  Mark-level top-$k$, nucleus, and per-family operations likewise fall
outside Proposition~\ref{prop:app-compose-refinement}; each would require its
own encoding-dependent conditions.
\hfill$\square$

\subsection{Finite-horizon terminal tilt}
\label{sec:app-compose-terminal-proof}

\begin{composeappendixlemma}[Backward normalization]
\label{lem:app-compose-backward-normalization}
Under A4, if $h_b(x)>0$, the row
$P_b^g(y\mid x)=P_b(y\mid x)h_{b-1}(y)/h_b(x)$ is row stochastic.
\end{composeappendixlemma}

\paragraph{Proof.}
Nonnegativity is immediate, and
\[
  \sum_yP_b^g(y\mid x)
  =
  \frac{\sum_yP_b(y\mid x)h_{b-1}(y)}{h_b(x)}
  =
  \frac{h_b(x)}{h_b(x)}
  =
  1.
\]
\hfill$\square$

\begin{composeappendixproposition}[Exact terminal tilt]
\label{prop:app-compose-terminal-tilt}
Fix an initial state $x_0$, horizon $K$, and nonnegative desirability $g$ with
$0<h_K(x_0)<\infty$.  The controlled chain defined by
Equation~\eqref{eq:finite-doob} has terminal law
\[
  \Pr^g(X_K=y\mid X_0=x_0)
  =
  \frac{\Pr(X_K=y\mid X_0=x_0)g(y)}{h_K(x_0)}.
\]
\end{composeappendixproposition}

\paragraph{Proof.}
For a base-positive path $x_{0:K}$ that ends at a state with $g(x_K)>0$,
every backward value encountered along that path is positive.  Its controlled
probability is therefore
\begin{align}
  \prod_{j=0}^{K-1}
  P_{K-j}^g(x_{j+1}\mid x_j)
  &=
  \prod_{j=0}^{K-1}
  P_{K-j}(x_{j+1}\mid x_j)
  \frac{h_{K-j-1}(x_{j+1})}{h_{K-j}(x_j)}
  \notag\\
  &=
  \left\{
  \prod_{j=0}^{K-1}
  P_{K-j}(x_{j+1}\mid x_j)
  \right\}
  \frac{h_0(x_K)}{h_K(x_0)}
  \notag\\
  &=
  \Pr(x_{1:K}\mid x_0)
  \frac{g(x_K)}{h_K(x_0)}.
  \label{eq:app-compose-path-telescope}
\end{align}
Paths ending where $g=0$ receive zero controlled mass.  Summing
Equation~\eqref{eq:app-compose-path-telescope} over all paths ending at $y$
gives the stated terminal law.  Summing again over $y$ shows that
$h_K(x_0)$ is its normalizer.  The identity holds at the declared horizon;
it does not identify an arbitrary early-stopped marginal with the terminal
tilt.
\hfill$\square$

\subsection{Support, equality, and zero backward values}
\label{sec:app-compose-support-proof}

\begin{composeappendixproposition}[Controlled support and unreachable rows]
\label{prop:app-compose-support}
For every $(b,x)$ with $h_b(x)>0$,
\[
  \operatorname{supp}P_b^g(x,\cdot)
  =
  \{y:P_b(y\mid x)>0,\ h_{b-1}(y)>0\}
  \subseteq
  \operatorname{supp}P_b(x,\cdot).
\]
The two supports are equal if and only if every base successor has positive
next-state value.  If $h_b(x)=0$, no positive-base-probability length-$b$
continuation from $x$ ends where $g>0$; the controlled row is undefined.
\end{composeappendixproposition}

\paragraph{Proof.}
On a defined row the denominator $h_b(x)$ is strictly positive, so
$P_b^g(y\mid x)>0$ holds exactly when both factors
$P_b(y\mid x)$ and $h_{b-1}(y)$ are positive.  This proves the support
identity, the subset relation, and the equality condition.

If $h_b(x)=0$, then
\[
  0=h_b(x)=\E[g(X_b)\mid X_0=x].
\]
Because $g(X_b)$ is nonnegative, it is zero almost surely under the base
continuation.  Equivalently, a positive-probability path to a state with
$g>0$ would make the expectation positive, which is a contradiction.
Equation~\eqref{eq:finite-doob} would divide zero successor weights by a zero
normalizer.  Adding an epsilon would define a different controller and may
assign mass without a conditioning interpretation, so the row must instead
be reported unreachable.
\hfill$\square$

\subsection{Dynamic continuation}
\label{sec:app-compose-dynamic-proof}

\begin{composeappendixproposition}[Exact retargeted continuation]
\label{prop:app-compose-dynamic}
Assume A4--A5.  At a switch state $x$ with $b$ edits remaining, replace the
old desirability by $g'$.  If $h_b^{g'}(x)>0$, recomputing the backward values
from $g'$ and applying their Doob rows produces
\[
  \Pr^{g'}(X_b=y\mid X_0=x)
  =
  \frac{\Pr(X_b=y\mid X_0=x)g'(y)}
       {h_b^{g'}(x)}.
\]
This is the exact retargeted continuation from $(b,x)$.
\end{composeappendixproposition}

\paragraph{Proof.}
Condition on the realized switch state and all context included in A5.  The
future base law is then the budget-indexed Markov chain with kernels
$P_b,\ldots,P_1$, independently of how the prefix reached $x$.
Proposition~\ref{prop:app-compose-terminal-tilt}, applied with initial state
$x$, horizon $b$, and terminal desirability $g'$, gives the displayed law.
The old objective affects which switch states were likely before
conditioning, but it does not enter the exact future transform once $x$ and
the sufficient context are fixed.
\hfill$\square$

The proposition does not say that retargeting reproduces the counterfactual
law obtained by using $g'$ from the original source.  Unconditionally, a
random switch state produces a mixture of state-conditional tilted
continuations.  Likewise, if a constraint depends on an unrecorded prefix
statistic, A5 fails; exact continuation then requires a killed process or an
augmented state.  At molecular scale, replacing the exact $h_b$ by a learned
$h_\phi$ makes the controller approximate even though its transitions remain
inside the base successor support.
